\documentclass[a4paper,12pt]{article}
\usepackage[utf8]{inputenc}
\usepackage{geometry}
\usepackage{graphicx}
\usepackage{listings}
\usepackage{xcolor}
\usepackage{float}
\usepackage{hyperref}

% Code highlighting configuration
\definecolor{codegreen}{rgb}{0,0.6,0}
\definecolor{codegray}{rgb}{0.5,0.5,0.5}
\definecolor{backcolour}{rgb}{0.95,0.95,0.92}

\lstdefinestyle{mystyle}{
    backgroundcolor=\color{backcolour},   
    commentstyle=\color{codegreen},
    keywordstyle=\color{magenta},
    basicstyle=\ttfamily\footnotesize,
    breakatwhitespace=false,         
    breaklines=true,                 
    captionpos=b,                    
    keepspaces=true,                 
    numbers=left,                    
    numbersep=5pt,                  
    showspaces=false,                
    showtabs=false,                  
    tabsize=2,
    frame=single
}
\lstset{style=mystyle}

\title{Practical Work 7: Sudoku on Google App Engine}
\author{Group Name \\ Member 1 \\ Member 2}
\date{\today}

\begin{document}

\maketitle

\section{Introduction}
Google App Engine (GAE) is a Platform as a Service (PaaS) offering that allows developers to build and run applications on Google's infrastructure. This practical work demonstrates how to build a simple Sudoku web application using Python (Flask) and prepare it for deployment on GAE.

\section{Project Structure}
The project follows the standard structure required by the GAE Python runtime environment:

\begin{itemize}
    \item \texttt{app.yaml}: The configuration file that defines the runtime environment and handlers.
    \item \texttt{main.py}: The core application logic utilizing the Flask micro-framework.
    \item \texttt{requirements.txt}: Lists dependencies (Flask, Gunicorn) to be installed by GAE.
    \item \texttt{templates/index.html}: The frontend HTML template for rendering the Sudoku board.
\end{itemize}

\section{Implementation Details}

\subsection{Configuration (app.yaml)}
This snippet defines the Python 3.9 runtime and the entry point for the application using Gunicorn.

\begin{lstlisting}[language=bash]
runtime: python39
entrypoint: gunicorn -b :$PORT main:app

handlers:
  - url: /.*
    script: auto
\end{lstlisting}

\subsection{Application Logic (main.py)}
We use a list of lists to represent the Sudoku board. The \texttt{index} function handles both GET requests (displaying the board) and POST requests (validating the user's solution).

\begin{lstlisting}[language=python]
from flask import Flask, render_template, request
app = Flask(__name__)

# Simplified Board Logic
@app.route('/', methods=['GET', 'POST'])
def index():
    board = PUZZLE_BOARD
    message = ""
    if request.method == 'POST':
        # Logic to validate user input against SOLVED_BOARD
        # Sets message to "Congratulations" or "Incorrect"
        pass 
    return render_template('index.html', board=board, message=message)
\end{lstlisting}

\subsection{Frontend Interface (templates/index.html)}
The HTML template utilizes Jinja2 templating engine to render the 9x9 grid dynamically. It uses nested loops to create input fields for each cell and submits the user's data via a POST form.

\begin{lstlisting}[language=html]
<!DOCTYPE html>
<html>
<body>
    <h1>Sudoku on Google App Engine</h1>
    <form method="POST">
        <table>
            {% for r in range(9) %}
            <tr>
                {% for c in range(9) %}
                <td>
                    <input type="text" name="cell-{{r}}-{{c}}" 
                           value="{{ board[r][c] if board[r][c] != 0 else '' }}" 
                           maxlength="1">
                </td>
                {% endfor %}
            </tr>
            {% endfor %}
        </table>
        <button type="submit" class="btn">Check Solution</button>
    </form>
    <div class="msg">{{ message }}</div>
</body>
</html>
\end{lstlisting}

\section{Setup and Execution}

Since modern Linux distributions adopt \texttt{PEP 668}, we cannot install Python packages system-wide. Therefore, we set up a virtual environment to run the application locally.

\subsection{Environment Setup}
We executed the following commands to prepare the environment and install dependencies:

\begin{lstlisting}[language=bash]
# 1. Create a virtual environment
sudo apt install python3-venv -y
python3 -m venv venv

# 2. Activate the environment
source venv/bin/activate

# 3. Install required libraries
pip install flask gunicorn
\end{lstlisting}

\subsection{Running the Application}
After the setup, we started the Flask development server:

\begin{lstlisting}[language=bash]
# Run the application
python3 main.py
\end{lstlisting}

The terminal output confirmed that the server was running on \texttt{http://127.0.0.1:8080}.

\subsection{Results}
We verified the gameplay by accessing the local web server. Figure \ref{fig:sudoku} displays the Sudoku game interface where users can input numbers and check their solutions.

\begin{figure}[H]
    \centering
    % Đổi tên file ảnh chụp màn hình của bạn thành sudoku_result.png
    \includegraphics[width=0.85\linewidth]{image1.png}
    \caption{Sudoku Game Interface running on Local Development Server}
    \label{fig:sudoku}
\end{figure}

\section{Conclusion}
In this practical work, we successfully structured a Python web application for the Google App Engine platform. By using Flask, we separated the logic (Python) from the presentation (HTML templates). 
\begin{itemize}
    \item The \texttt{app.yaml} file is crucial for GAE to understand how to serve the application.
    \item The PaaS model allows us to focus on the code (\texttt{main.py}) without worrying about the underlying infrastructure setup.
\end{itemize}

\end{document}
